\section{Related Work}
\label{sec:related}

% Full 21-system prior-art table and per-system descriptions relocated to Appendix~G.

We organize related work into four categories and position TSCG within the design space.

\paragraph{Compression-based methods.}
LLMLingua~\cite{jiang2023llmlingua} and LLMLingua-2~\cite{pan2024llmlingua2} achieve up to 20$\times$ compression on natural prose through perplexity-based token importance scoring, but require GPU model inference, produce non-deterministic output, and---critically---are ineffective on structured content (80\% accuracy on tool schemas vs.\ TSCG's 93.3\%; the compound pipeline yields 0\%, \S\ref{sec:bfcl-integration}).
Selective Context~\cite{li2023compressing} uses self-information thresholds but lacks tokenizer and positional awareness.
CompactPrompt~\cite{compactprompt2025} is the only other LLM-free heuristic, achieving 1.44$\times$ via n-gram abbreviation---versus TSCG's 3.5$\times$ on structured content---limited by lexical-level operation.

\paragraph{Search-based methods.}
DSPy~\cite{khattab2023dspy} compiles declarative LLM calls into self-improving pipelines via bootstrapped demonstrations.
SAMMO~\cite{schlegel2024sammo} shares TSCG's ``compile-time'' framing but requires an LLM for search-based mutations and produces non-deterministic output.
OPRO~\cite{yang2024large}, APE~\cite{zhou2023large}, and EvoPrompt~\cite{chen2024evoprompt} all require iterative model interaction.
These systems optimize prompt \emph{content} (what to say); TSCG optimizes prompt \emph{structure} (how to represent it).

\paragraph{Gradient and template methods.}
TextGrad~\cite{yuksekgonul2024textgrad} introduces textual gradients for prompt refinement but requires model access.
LangGPT~\cite{wang2024langgpt} provides structured templates without formal optimization.
Neither addresses tokenizer alignment or causal attention grounding.

\paragraph{Tool-calling benchmarks.}
BFCL~\cite{patil2024gorilla} evaluates function-calling accuracy but always with full, uncompressed schemas.
ToolBench~\cite{qin2023toolllm} tests multi-step API planning across 16{,}000+ APIs without considering schema token overhead.
The prompt compression survey~\cite{li2024prompt} evaluates no system on function-calling schemas.
TAB fills this gap: the first benchmark measuring tool-schema compression effects on LLM tool-use performance.

\subsection{Design-Space Comparison}

Table~\ref{tab:design-space} summarizes TSCG's positioning along six axes.

\begin{table}[h]
\centering
\small
\caption{Design-space comparison of representative prompt optimization systems.}
\label{tab:design-space}
\begin{tabular}{lcccccc}
\toprule
\textbf{System} & \textbf{Determ.} & \textbf{Zero Dep.} & \textbf{Theory} & \textbf{Schema} & \textbf{Compress.} & \textbf{Speed} \\
\midrule
LLMLingua-2 & \xmark & \xmark & \xmark & \xmark & 20$\times$ & 42\,s \\
DSPy & \xmark & \xmark & \xmark & \xmark & --- & min \\
TextGrad & \xmark & \xmark & \xmark & \xmark & --- & min \\
Sel.\ Context & \xmark & \xmark & \xmark & \xmark & 5$\times$ & ms \\
LangGPT & \cmark & \cmark & \xmark & \xmark & --- & ms \\
\textbf{TSCG} & \cmark & \cmark & \cmark & \cmark & \textbf{3.5$\times$} & \textbf{$<$1\,ms} \\
\bottomrule
\end{tabular}
\end{table}

\noindent No existing system combines deterministic output, zero dependencies, theoretical grounding in attention mechanics, and specialization for tool schemas.
TSCG's niche---structured prompt content in agentic systems---is underserved: the NAACL 2025 compression survey evaluates no system on function-calling schemas.
On natural prose, LLMLingua-2 dominates (20$\times$ vs.\ TSCG's 1.07$\times$); on prompt content optimization, DSPy achieves complementary gains.
TSCG is the strongest system for a specific, growing problem---not a general-purpose replacement.
